\chapter*{Glossary}
\addcontentsline{toc}{chapter}{Glossary}

% ---------------------------------------------------------------
% Glossary of domain terms used in the book. Starter list covers
% terms that recur across chapters. Expand as needed; write each
% definition in the book's voice (plain language, units where relevant).
% ---------------------------------------------------------------

\begin{description}
  \item[Canary deployment]
        A release strategy where a new model version
        serves a small fraction of traffic while metrics and safety
        signals are monitored, before ramping to full rollout.

  \item[Context window]
        The maximum number of input tokens a model can
        process in a single inference call.

  \item[Drift]
        A change over time in the distribution of model
        inputs, outputs, or quality signals that threatens the
        assumptions under which a deployment was approved.

  \item[Embedding]
        A dense vector representation of text used for
        retrieval, clustering, or similarity search.

  \item[Grounding / groundedness]
        The property that a generated claim is supported
        by retrieved evidence; measured by evaluator models or rubrics
        in production.

  \item[Hallucination]
        A plausible-sounding but factually incorrect or
        unsupported generation; a first-class failure mode for LLM
        systems.

  \item[LLM-as-judge]
        An evaluation pattern in which a model scores
        another model's outputs; requires bias mitigations to be
        trustworthy as a gate.

  \item[LoRA (Low-Rank Adaptation)]
        A parameter-efficient fine-tuning method that
        trains a small low-rank update to the base model's weights
        instead of all parameters.

  \item[Observability]
        The ability to diagnose a system's behavior from
        its external outputs---traces, metrics, and logs---without
        having to re-run it.

  \item[Prompt template]
        A versioned, parameterized instruction used to
        elicit a specific model behavior; treated as a released
        artifact.

  \item[Quantization]
        Reducing the numerical precision of model weights
        or activations (e.g., FP16 to INT8) to lower memory and
        compute cost, typically at some quality trade-off.

  \item[Retrieval-Augmented Generation (RAG)]
        An inference pattern that combines a retrieval
        step over an external corpus with a generation step to ground
        responses in authoritative sources.

  \item[SLO (Service-Level Objective)]
        A numeric target for a user-visible service
        property (latency, availability, quality) that a team commits
        to meet over a defined window.

  \item[TTFT (Time-To-First-Token)]
        The latency between request receipt and the first
        generated token; a primary latency SLO for streamed responses.

  \item[Vector store]
        A database that indexes embeddings for efficient
        nearest-neighbor retrieval, typically with metadata filters.

  % Add more terms here as needed.
\end{description}
